\documentclass[11pt,reqno]{amsart}
\usepackage[localtheorems]{raeez-math-template}

\usepackage{array}
\usepackage{tikz-cd}

\newtheorem{theorem}{Theorem}[section]
\newtheorem{proposition}[theorem]{Proposition}
\newtheorem{lemma}[theorem]{Lemma}
\newtheorem{corollary}[theorem]{Corollary}
\newtheorem{conditionaltheorem}[theorem]{Conditional Theorem}
\theoremstyle{definition}
\newtheorem{definition}[theorem]{Definition}
\newtheorem{construction}[theorem]{Construction}
\newtheorem{example}[theorem]{Example}
\newtheorem{problem}[theorem]{Open Problem}
\theoremstyle{remark}
\newtheorem{remark}[theorem]{Remark}

\providecommand{\Ran}{\mathrm{Ran}}
\providecommand{\Sym}{\mathrm{Sym}}
\providecommand{\FactD}{\mathrm{FactD}}
\providecommand{\DMod}{\mathrm{DMod}}
\providecommand{\Alg}{\mathrm{Alg}}
\providecommand{\ComCoalg}{\mathrm{ComCoalg}}
\providecommand{\Fin}{\mathrm{Fin}}
\providecommand{\sdim}{\mathrm{sdim}}
\providecommand{\CE}{\mathrm{CE}}
\newcommand{\Barc}{\operatorname{Bar}}
\providecommand{\den}{\mathrm{den}}
\providecommand{\id}{\mathrm{id}}
\providecommand{\unit}{\mathbf{1}}
\providecommand{\Hd}{\boxtimes_{\mathrm{H}}}
\providecommand{\Cy}{\boxtimes_{\mathrm{C}}}

\title[Species, locality, and bar denominators]
{Duoidal linear species, factorization locality, and bar denominator products}

\author{Raeez Lorgat}

\date{September 2026}

\begin{document}

\begin{abstract}
The Cauchy and Hadamard products on linear species admit canonical
braided duoidal structures in both orientations.  Their interchange
maps are diagonal inclusion and off-diagonal-killing projection, and
\(\iota\pi\) is the idempotent onto the common-decomposition
summand.  These structures are globally non-strong.  Independently, a
factorization right \(D\)-module on the reduced Ran space carries
collision descent and coherent factorization equivalences.  An ambient
augmented \(E_1\)-structure is factorization-local precisely when its
entire operadic action lifts through the factorization category; the
usual disjointness-locus multiplication square is only a necessary
binary consequence.

For a positively root-graded Lie superalgebra satisfying a
finite-factorization condition, the completed total Euler
supercharacter of the normalized algebraic bar of
\(U(\mathfrak n_+)\) is the signed root product.  Under the precise
hypotheses of a superdenominator theorem this is its product factor.
For the Gritsenko--Nikulin algebra, a shifted formal Fourier map followed
by a normally convergent cusp realization gives
\(D_5(2Z)=64^{-1}\Delta_5(2Z)\).  Weightwise duality with Goncharova's
continuous cohomology gives the two positive-Witt homology classes in
weights \((3n^2\pm n)/2\).  The Motzkin--Riordan relation is proved
independently by its generating functions.  A comparison between the
factorization bar and the root-graded algebraic bar is stated
conditionally; no such comparison is used in the proved results.
\end{abstract}

\maketitle

\section{The comparison problem}

Dividing a set of insertions into two groups gives an ordered
decomposition.  Two operations can use the same decomposition or choose
their decompositions independently.  For a singleton there are two
ordered decompositions, so these choices give two and four possibilities,
respectively.  Inclusion of the common choices and projection onto them
give a retraction.  They do not give an invertible exchange.

Linear species assign vector spaces to finite sets and transport them
along bijections.  The Cauchy product distributes a set between its two
inputs.  The Hadamard product evaluates both inputs on the whole set.
Their interchange maps implement the preceding inclusion and projection.
The resulting two duoidal structures retain both products and their
coherent compatibility, with different units.  Noninvertibility is part
of this construction and leaves its bimonoids well defined.

When insertion points lie on a curve, they can also collide.  A
factorization right \(D\)-module specifies coefficients on configurations,
their restriction to collision diagonals, and their decomposition away
from those diagonals.  An ordered product on it must preserve all these
maps and their higher coherences.  Such a product defines an internal
simplicial bar.  Its realization in the factorization category requires
the colimit hypotheses of Theorem~\ref{thm:realized-bar}.

A second calculation starts from root-graded Lie superalgebras.
The normalized bar of their enveloping algebras has an Euler
supercharacter determined by the even and odd root spaces.
Finite factorization of weights makes each coefficient a finite
calculation.  The resulting product has a precise Igusa specialization,
including a Weyl shift and an analytic domain of convergence.
To obtain that character from a factorization bar, one still needs a map
between the two chain constructions.  Problem~\ref{prob:carrier-change}
specifies its coefficients and operations, and
Theorem~\ref{thm:conditional-comparison} gives the consequence once those
data are constructed.

\section{Conventions}

Throughout, \(k\) is a field of characteristic zero.  The groupoid
\(\Fin^{\simeq}\) contains all finite sets, including
\(\varnothing\), and species are covariant under bijections.  By
contrast, \(\Fin^{\mathrm{surj}}\) denotes the category of nonempty
finite sets and surjections.  It indexes the reduced Ran space; there
is no empty Ran configuration in this convention.

Super vector spaces are \(\mathbb Z/2\)-graded and their morphisms
preserve intrinsic parity.  Mapping complexes, dg objects, and
\(A_\infty\)-algebras use cohomological degrees: the differential has
degree \(+1\), \((A[d])^p=A^{p+d}\), and \(m_n\) has degree \(2-n\).
Normalized simplicial bars and Chevalley--Eilenberg chains carry a
separate nonnegative homological degree \(r\), whose differential lowers
\(r\) by one.  In a total Euler supercharacter a class of this
homological degree and intrinsic parity \(\epsilon\) has sign
\((-1)^{r+\epsilon}\).  These two grading conventions are not
identified.  When a cohomological degree \(p\) and a bar degree \(r\)
occur together, the cohomological totalization has degree \(p-r\).

For the geometric statements, \(X\) is a smooth separated complex
algebraic curve.  All categories of \(D\)-modules are stable
\(\infty\)-categories of right \(D\)-modules.  Extraordinary pullback is
denoted by \((-)^!\), and every collision, factorization, and refinement
equivalence is understood homotopy coherently.

\section{The two duoidal structures on linear species}

Let
\[
 \operatorname{Sp}_k
 =\operatorname{Fun}(\Fin^{\simeq},\operatorname{Vect}_k).
\]
For \(F,G\in\operatorname{Sp}_k\), define the Cauchy and Hadamard
products by
\[
 (F\Cy G)(I)
 =\bigoplus_{I=I_1\sqcup I_2}F(I_1)\otimes G(I_2),
 \qquad
 (F\Hd G)(I)=F(I)\otimes G(I).
\]
The decompositions in the first sum are ordered.  Write \(\mathbb U\)
for the Cauchy unit and \(\mathbb E\) for the Hadamard unit:
\[
 \mathbb U(\varnothing)=k,\qquad
 \mathbb U(I)=0\quad(I\ne\varnothing),\qquad
 \mathbb E(I)=k.
\]

For two ordered decompositions
\(d=(I_1,I_2)\) and \(e=(J_1,J_2)\) of \(I\), put
\[
 T_{d,e}
 =F_1(I_1)\otimes F_2(I_2)\otimes
  G_1(J_1)\otimes G_2(J_2).
\]
After the evident symmetry of tensor factors,
\[
 \begin{split}
 L_I&:=((F_1\Hd G_1)\Cy(F_2\Hd G_2))(I)
      \cong\bigoplus_dT_{d,d},\\
 R_I&:=((F_1\Cy F_2)\Hd(G_1\Cy G_2))(I)
      \cong\bigoplus_{d,e}T_{d,e}.
 \end{split}
\]
Diagonal inclusion and projection define natural transformations
\[
 \iota:L\longrightarrow R,\qquad
 \pi:R\longrightarrow L.
\]

\begin{theorem}[Cauchy--Hadamard duoidality]
\label{thm:species-duoidal}
The transformations \(\iota\) and \(\pi\) define braided duoidal
structures in the two respective orientations
\[
 (\operatorname{Sp}_k,\Cy,\mathbb U,\Hd,\mathbb E;\iota),
 \qquad
 (\operatorname{Sp}_k,\Hd,\mathbb E,\Cy,\mathbb U;\pi).
\]
For the \(\iota\)-orientation the three unit morphisms are
\[
 \mathbb U\longrightarrow\mathbb U\Hd\mathbb U,\qquad
 \mathbb E\Cy\mathbb E\longrightarrow\mathbb E,\qquad
 \mathbb U\longrightarrow\mathbb E.
\]
For the \(\pi\)-orientation they are
\[
 \mathbb E\longrightarrow\mathbb E\Cy\mathbb E,\qquad
 \mathbb U\Hd\mathbb U\longrightarrow\mathbb U,\qquad
 \mathbb E\longrightarrow\mathbb U.
\]
Moreover,
\[
 \pi\iota=\id_L,\qquad e:=\iota\pi,\qquad e^2=e,
\]
and \(e\) has image \(\bigoplus_dT_{d,d}\).
\end{theorem}

\begin{proof}
The map \(\iota_I\) includes the summands indexed by
\(d\mapsto(d,d)\), while \(\pi_I\) is the identity on those summands
and zero on every \(T_{d,e}\) with \(d\ne e\).  Hence
\(\pi_I\iota_I=\id\), and \(\iota_I\pi_I\) is the displayed
idempotent.

The unit maps contain information not visible in this retraction.  The
map \(\mathbb E\Cy\mathbb E\to\mathbb E\) is the fold over all ordered
decompositions, and
\(\mathbb E\to\mathbb E\Cy\mathbb E\) is the biproduct diagonal.
The maps \(\mathbb U\to\mathbb E\) and
\(\mathbb E\to\mathbb U\) are the identity at \(\varnothing\) and zero
elsewhere.  The two maps involving
\(\mathbb U\Hd\mathbb U\) are the canonical unit identifications at
\(\varnothing\) and the unique maps between zero spaces elsewhere.

Every associativity diagram decomposes into matrix entries indexed by
ordered decompositions into three blocks.  In the \(\iota\)-orientation
both composites retain exactly those entries whose decomposition labels
agree; in the \(\pi\)-orientation both composites kill exactly the
complementary entries.  The unit and braided diagrams reduce in the
same way to the fold, diagonal, identity, or zero maps just described.
This proves the coherence axioms.  For vector species these two
orientations are also the braided \(2\)-monoidal structures of
\cite{AM10}*{Definitions 6.1 and 6.3, equations (8.73)--(8.74),
Propositions 8.58 and 8.68};
the three-map unit datum is the one in
\cite{Tor25}*{Definition 2.1}.
\end{proof}

\begin{proposition}[The common-decomposition summand]
\label{prop:common-summand}
For fixed \(F_1,F_2,G_1,G_2\) and \(I\), the map \(\iota_I\) is an
isomorphism if and only if every off-diagonal summand \(T_{d,e}\),
\(d\ne e\), is zero.  This is an objectwise condition; it does not
define a unital monoidal subcategory containing the inherited units.
\end{proposition}

\begin{proof}
The target is the direct sum of the image of \(\iota_I\) and all
off-diagonal summands.  Thus \(\iota_I\) is surjective exactly when the
latter vanish.
\end{proof}

\begin{corollary}[Failure of strongness]
\label{cor:not-strong}
If \(k\ne0\), neither duoidal structure in
Theorem~\ref{thm:species-duoidal} is strong.  There is no full
subcategory containing both inherited units on which all interchange
and unit morphisms become isomorphisms.
\end{corollary}

\begin{proof}
Take \(I=\{*\}\) and all four inputs equal to \(\mathbb E\).  There are
two ordered decompositions of \(I\), and therefore
\[
 L_I\cong k^2,\qquad R_I\cong k^4.
\]
In compatible coordinates,
\[
 \iota(a,b)=(a,0,0,b),\qquad
 \pi(a,b,c,d)=(a,d),\qquad
 \iota\pi=\operatorname{diag}(1,0,0,1).
\]
The interchange is not invertible.  Independently,
\(\mathbb U(\{*\})=0\) and \(\mathbb E(\{*\})=k\), so neither
cross-unit map is invertible.  Any subcategory containing the inherited
units contains this obstruction.
\end{proof}

\begin{remark}
Bimonoids are defined for a duoidal category without requiring strong
interchange.  The idempotent \(e\) records a useful summand, but
invertibility is neither part of duoidality nor needed for the
definition of a bimonoid.
\end{remark}

\section{Factorization locality on the reduced Ran space}

Species record decomposition labels but have no maps for collisions of
points on \(X\).  Collision diagonals supply these maps.  The
factorization condition then relates their restrictions to products on
disjoint configurations.

For a surjection \(q:I\twoheadrightarrow J\), let
\[
 \Delta_q:X^J\longrightarrow X^I
\]
be the collision diagonal.  The reduced Ran category of right
\(D\)-modules is
\[
 \DMod^r(\Ran X)
 =\lim_{I\in\Fin^{\mathrm{surj}}}\DMod^r(X^I),
\]
where the transition attached to \(q\) is
\[
 \Delta_q^!:\DMod^r(X^I)\longrightarrow\DMod^r(X^J).
\]
Thus a Ran object consists of objects \(A_I\in\DMod^r(X^I)\) and
equivalences
\[
 \theta_q:\Delta_q^!A_I\xrightarrow{\sim}A_J,
\]
with coherent identity and composition data
\[
 \theta_{r\circ q}\simeq
 \theta_r\circ\Delta_r^!(\theta_q)
\]
for composable surjections.

For \(\pi:I\twoheadrightarrow J\), put \(I_j=\pi^{-1}(j)\) and
\[
 U(\pi)=
 \{(x_i)\in X^I:x_i\ne x_{i'}
       \text{ whenever }\pi(i)\ne\pi(i')\},
 \qquad
 j_\pi:U(\pi)\hookrightarrow X^I.
\]

\begin{definition}[Factorization right \(D\)-module]
\label{def:factorization}
A factorization right \(D\)-module is a Ran object \(A\) together with
equivalences
\[
 \phi_\pi:
 j_\pi^!A_I\xrightarrow{\sim}
 j_\pi^!\left(\boxtimes_{j\in J}A_{I_j}\right)
\]
for every \(\pi\), compatibly with collision maps, bijections,
refinements of surjections, and every higher refinement simplex.
Morphisms respect all \(\theta_q,\phi_\pi\), and their coherences.
\end{definition}

We write \(\FactD(X)\) for the resulting factorization subcategory,
equivalently the Francis--Gaitsgory category
\[
 \FactD(X)
 =
 \ComCoalg^{\mathrm{ch}}_{\mathrm{Fact}}
   \bigl(\DMod^r(\Ran X)\bigr).
\]
It includes in the category of chiral commutative coalgebras, which in
turn has a forgetful functor to \(\DMod^r(\Ran X)\).  The composite is
not a full embedding of \(\FactD(X)\) into underlying Ran
\(D\)-modules.  The reduced Ran limit,
chiral convolution, and factorization condition are given in
\cite{FG12}*{Section 2.1.1, Definition 2.1.2, Lemma 2.3.4,
Definitions 2.4.5 and
2.4.7}; the coherent-system interpretation is
\cite{FG12}*{Remark 2.4.8}.

\subsection{The pointwise tensor}

For right \(D\)-modules define
\[
 (A\otimes^!B)_I=A_I\otimes^!B_I.
\]
The collision equivalence is the composite
\[
 \Delta_q^!(A_I\otimes^!B_I)
 \simeq
 \Delta_q^!A_I\otimes^!\Delta_q^!B_I
 \xrightarrow{\theta_q^A\otimes\theta_q^B}
 A_J\otimes^!B_J.
\]
The factorization equivalence is
\[
\begin{aligned}
 j_\pi^!(A_I\otimes^!B_I)
 &\simeq j_\pi^!A_I\otimes^!j_\pi^!B_I\\
 &\xrightarrow{\phi_\pi^A\otimes\phi_\pi^B}
 j_\pi^!\!\left(\boxtimes_jA_{I_j}\right)
 \otimes^!
 j_\pi^!\!\left(\boxtimes_jB_{I_j}\right)\\
 &\xrightarrow{\mathrm{ex}}
 j_\pi^!\!\left(
   \boxtimes_j(A_{I_j}\otimes^!B_{I_j})\right).
\end{aligned}
\]
The exchange map and its coherence are part of this construction.
Symmetric monoidality of extraordinary pullback and external-product
exchange supplies the identity, composition, permutation, and
refinement coherences.  They give the symmetric monoidal category
\[
 (\FactD(X),\otimes^!,\boldsymbol\omega_{\Ran}),
 \qquad
 (\boldsymbol\omega_{\Ran})_I=\omega_{X^I},
\]
and a strong symmetric monoidal forgetful functor
\[
 U:\FactD(X)\longrightarrow\DMod^r(\Ran X).
\]
The convention remains reduced: there is no component
\(A_\varnothing\).  An augmentation is a morphism
\(A\to\boldsymbol\omega_{\Ran}\) for the pointwise tensor and does not
adjoin an empty configuration to the chiral coalgebra.

\begin{definition}[Ordered factorization algebra]
\label{def:ordered-factorization}
An ordered factorization algebra is an object
\[
 A\in
 \Alg^{\mathrm{aug}}_{E_1}\bigl(\FactD(X),\otimes^!\bigr).
\]
\end{definition}

\begin{proposition}[Locality as an operadic lift]
\label{prop:locality-lift}
Let \(A\in\FactD(X)\), and let \(\alpha\) be an augmented
\(E_1\)-algebra structure on \(U(A)\).  A factorization-local structure
on \(\alpha\) is precisely a lift of maps of \(\infty\)-operads
\[
\begin{tikzcd}[column sep=4.5em,row sep=3em]
 & \FactD(X)^\otimes
     \arrow[d,"U^\otimes"]\\
 (E_1^{\mathrm{aug}})^\otimes
     \arrow[r,"\alpha"']
     \arrow[ur,dashed,"\widetilde\alpha"]
 & \DMod^r(\Ran X)^\otimes .
\end{tikzcd}
\]
Equivalently, multiplication, unit, augmentation, every higher
operation, and all coherent homotopies lift to the mapping spaces of
\(\FactD(X)\), with the algebra color sent to the specified object
\(A\).
\end{proposition}

\begin{proof}
An augmented \(E_1\)-algebra is a map from the augmented associative
\(\infty\)-operad into the monoidal \(\infty\)-category.  Lifting that
map along \(U^\otimes\) is exactly the definition of an augmented
\(E_1\)-algebra structure on \(A\) whose underlying structure is
\(\alpha\).
\end{proof}

The familiar multiplication square is a necessary arity-two
consequence.  For an ordered bipartition \(d=(I_1,I_2)\) with
\(I_1,I_2\ne\varnothing\), equivalently a surjection
\(I\twoheadrightarrow\{1,2\}\), define
\[
\begin{aligned}
 \lambda_d^{(2)}:
 j_d^!(A_I\otimes^!A_I)
 &\simeq j_d^!A_I\otimes^!j_d^!A_I\\
 &\xrightarrow{\phi_d\otimes\phi_d}
 j_d^!(A_{I_1}\boxtimes A_{I_2})
 \otimes^!
 j_d^!(A_{I_1}\boxtimes A_{I_2})\\
 &\xrightarrow{\mathrm{ex}}
 j_d^!\!\left(
 (A_{I_1}\otimes^!A_{I_1})\boxtimes
 (A_{I_2}\otimes^!A_{I_2})\right).
\end{aligned}
\]
Every lift in Proposition~\ref{prop:locality-lift} makes
\[
\begin{tikzcd}[column sep=4.2em,row sep=3em]
 j_d^!(A_I\otimes^!A_I)
   \arrow[r,"j_d^!m_{2,I}"]
   \arrow[d,"\lambda_d^{(2)}"']
 &
 j_d^!A_I
   \arrow[d,"\phi_d"]
 \\
 j_d^!\!\left(
 (A_{I_1}\otimes^!A_{I_1})\boxtimes
 (A_{I_2}\otimes^!A_{I_2})\right)
   \arrow[r,"j_d^!(m_{2,I_1}\boxtimes m_{2,I_2})"']
 &
 j_d^!(A_{I_1}\boxtimes A_{I_2})
\end{tikzcd}
\]
commute.  It also makes every collision square commute:
\[
\begin{tikzcd}[column sep=4em,row sep=3em]
 \Delta_q^!(A_I\otimes^!A_I)
   \arrow[r,"\Delta_q^!m_{2,I}"]
   \arrow[d,"\simeq"']
 &
 \Delta_q^!A_I
   \arrow[d,"\theta_q"]
 \\
 A_J\otimes^!A_J
   \arrow[r,"m_{2,J}"']
 &
 A_J .
\end{tikzcd}
\]
The left vertical map uses the monoidality of \(\Delta_q^!\), followed
by \(\theta_q\otimes\theta_q\).

After representing the \(E_1\)-structure by a cohomologically graded
\(A_\infty\)-algebra, analogous diagrams are required for every
\[
 m_n:A^{\otimes^!n}\longrightarrow A[2-n],
\]
as well as for the unit, augmentation, associahedral homotopies,
arbitrary multiblock surjections, refinements, and permutations.

\begin{example}[Binary compatibility is not sufficient]
\label{ex:binary-not-sufficient}
Assume \(\operatorname{char}k\ne2\).  In
\(\operatorname{Fun}(BC_2,\operatorname{Ch}_k)\), let
\[
 A^0=k1\oplus kx,\qquad A^{-1}=ky,
\]
with zero differential and square-zero augmentation ideal for \(m_2\).
Give \(A\) the strictly unital minimal \(A_\infty\)-operation
\[
 m_3(x,x,x)=y
\]
and no other nonzero higher operation.  Let the involution fix \(1,x\)
and send \(y\) to \(-y\).  Multiplication, unit, and augmentation are
equivariant, but \(m_3\) is not.  Thus binary equivariance does not
imply descent of an \(E_1\)-structure.
\end{example}

\begin{remark}[The species and Ran carriers]
\label{rem:species-ran}
The transformations \(\iota,\pi\) take values in one fixed additive
category.  The components of a Ran object lie in the varying categories
\(\DMod^r(X^I)\) and are joined by extraordinary pullback, open
restriction, and collision descent.  A geometric analogue of
\(\iota\) or \(\pi\) therefore requires mixed pull--push functors,
projection-formula maps, Beck--Chevalley compatibilities, collision
squares, units, and three-block refinement coherence.

There is also a first-arity discrepancy.  Reduced binary chiral
convolution has no singleton summand, since there is no surjection
\(\{*\}\twoheadrightarrow\{1,2\}\).  The Cauchy product of species has
two singleton summands arising from empty blocks.  Agreement of labels
on a fixed separated stratum consequently supplies neither a Ran
natural transformation nor a retraction.
\end{remark}

\section{The internal bar and its realization}

An augmented ordered product gives operations for composing adjacent
insertions and for removing either endpoint.  Their compatibility is
exactly the simplicial structure used by the bar construction.  Because
the product is factorization-local, these operations remain morphisms
of factorization objects before any realization is taken.

\begin{proposition}[The internal simplicial bar]
\label{prop:internal-bar}
Let
\(\mathcal F=(\FactD(X),\otimes^!,\boldsymbol\omega_{\Ran})\), and let
\(A\in\Alg^{\mathrm{aug}}_{E_1}(\mathcal F)\).  Then
\[
 \Barc^{\mathcal F}_r(\boldsymbol\omega_{\Ran},A,
 \boldsymbol\omega_{\Ran})
 =
 \boldsymbol\omega_{\Ran}\otimes^!
 A^{\otimes^!r}\otimes^!\boldsymbol\omega_{\Ran}
\]
defines a functor \(\Delta^{op}\to\mathcal F\).  The inner faces are
induced by multiplication, the two outer faces by the augmentation, and
the degeneracies by the unit.
\end{proposition}

\begin{proof}
Every simplicial operator is a composite of tensor products of
morphisms in \(\mathcal F\).  The simplicial identities are the coherent
unital associativity identities of the augmented \(E_1\)-algebra.
\end{proof}

\begin{conditionaltheorem}[Realized factorization bar]
\label{thm:realized-bar}
Assume the geometric realization of the simplicial object in
Proposition~\ref{prop:internal-bar} exists in \(\mathcal F\), and assume
the canonical map
\[
 \left|U\Barc^{\mathcal F}_\bullet
   (\boldsymbol\omega_{\Ran},A,\boldsymbol\omega_{\Ran})
 \right|_{\DMod^r(\Ran X)}
 \longrightarrow
 U\!\left(
 \left|\Barc^{\mathcal F}_\bullet
   (\boldsymbol\omega_{\Ran},A,\boldsymbol\omega_{\Ran})
 \right|_{\mathcal F}\right)
\]
is an equivalence.  Then the realized object is a factorization right
\(D\)-module whose underlying Ran object is the ambient realized bar.
If the ambient tensor preserves geometric realizations separately, that
underlying object computes
\[
 \boldsymbol\omega_{\Ran}
 \otimes^{\mathbb L}_{U(A)}
 \boldsymbol\omega_{\Ran}.
\]
\end{conditionaltheorem}

\begin{remark}[Normalization and deconcatenation]
\label{rem:deconcatenation}
Suppose normalized reduced chains and the necessary sums exist in a
chosen additive monoidal enhancement, and first suppose that \(A\) has
a strict augmented dg-associative representative.  With
\(\bar A=\operatorname{fib}(A\to\boldsymbol\omega_{\Ran})\), let
\(s\bar A\) denote a copy of \(\bar A\) placed in independent bar
homological degree one; \(s\) is not the cohomological shift \([1]\).
The reduced bar is the bigraded coalgebra
\[
 \overline B(A)=T^c(s\bar A)
 =\bigoplus_{r\ge0}(s\bar A)^{\otimes r}.
\]
Its internal differential raises cohomological degree and preserves
\(r\); its simplicial bar differential is the coderivation induced by
multiplication and lowers \(r\) by one.  Both have degree \(+1\) after
the totalization \(p-r\).
For a general cohomological \(A_\infty\)-representative, the full bar
coderivation instead has a component induced by every operation \(m_n\);
the component associated with \(m_n\) changes bidegree by
\((2-n,1-n)\) and hence again has total degree \(+1\).
Consequently
\[
 \Delta_{\mathrm{dec}}[a_1|\cdots|a_r]
 =\sum_{i=0}^r
 [a_1|\cdots|a_i]\otimes[a_{i+1}|\cdots|a_r]
\]
is a coassociative chain map in every word length.  To obtain this map
inside \(\FactD(X)\), normalization and the required sums must be
created by \(U\), tensor must preserve them, and the resulting map must
respect collision and factorization descent.  These colimits are not
automatic: the standard category \(\operatorname{Fact}(X)\) is not
presentable and lacks coproducts
\cite{FG12}*{Remark 2.4.9}.
\end{remark}

\section{The algebraic bar denominator}

For an enveloping algebra, antisymmetrization compares ordered bar words
with Chevalley--Eilenberg chains.  A positive root grading makes their
Euler supercharacter computable one weight at a time.  We first specify
the finiteness and sign conventions needed for this calculation.

In \(\operatorname{Ch}_{\ge0}(\operatorname{sVect}_k)\), write \(sV\)
for a super vector space \(V\) placed in homological degree one, with
its intrinsic parity unchanged.  Thus \(s\) is a homological suspension,
not the cohomological shift \([1]\).  On bidegrees
\((r,\epsilon)\) and \((r',\epsilon')\), the symmetry has sign
\((-1)^{rr'+\epsilon\epsilon'}\), while the Euler sign of one
bidegree is \((-1)^{r+\epsilon}\).

\begin{definition}[Finite-factorization monoid]
\label{def:finite-factorization}
A pointed commutative monoid \(Q_+\) has the
\emph{finite-factorization property} if, for every \(\beta\in Q_+\),
\[
 \coprod_{r\ge0}
 \left\{(\alpha_1,\ldots,\alpha_r)
 \in(Q_+\setminus\{0\})^r:
 \alpha_1+\cdots+\alpha_r=\beta\right\}
\]
is finite.
\end{definition}

This condition makes
\[
 \widehat{\mathbb Z[Q_+]}
 =\prod_{\beta\in Q_+}\mathbb Z X^\beta,
 \qquad
 (ab)_\beta=\sum_{\gamma+\delta=\beta}a_\gamma b_\delta,
\]
a ring.  It is stronger than finite-dimensionality of the homogeneous
pieces.  For example, the pointed monoid generated inside
\(\mathbb Z^2\) by \((1,j)\) and \((0,-j)\), \(j\ge1\), has infinitely
many decompositions
\[
 (1,0)=(1,j)+(0,-j).
\]

Let
\[
 \mathfrak n_+
 =\bigoplus_{\alpha\in Q_+\setminus\{0\}}\mathfrak g_\alpha
\]
be a \(Q_+\)-graded Lie superalgebra with finite-dimensional
homogeneous pieces.  Give \(k\) the trivial
\(U(\mathfrak n_+)\)-module structure and define
\[
 \mathcal B(\mathfrak n_+)
 :=
 N\Barc_\bullet(k,U(\mathfrak n_+),k)
 \in
 \operatorname{Ch}_{\ge0}
   (\operatorname{sVect}^{Q_+}_k).
\]
In simplicial homological degree \(r\),
\[
 \mathcal B_r(\mathfrak n_+)
 \cong\overline{U(\mathfrak n_+)}^{\otimes r}.
\]
Equivalently, its reduced dg-bar presentation is
\(T^c(s\overline{U(\mathfrak n_+)})\).
Normalization is essential: for the augmented algebra \(k\), the
normalized bar is zero in every positive degree, whereas the
unnormalized simplicial object has a copy of \(k\) in every degree.

Let
\[
 \widehat{\operatorname{Ch}}^{\mathrm{pf}}_{Q_+}
   (\operatorname{sVect}_k)
 :=
 \prod_{\beta\in Q_+}
   \operatorname{Perf}(\operatorname{sVect}_k)
\]
be the category of \(Q_+\)-graded supercomplexes that are bounded and
finite-dimensional in each weight.  Its completed convolution is
\[
 (C\widehat\otimes D)_\beta
 =\bigoplus_{\gamma+\delta=\beta}C_\gamma\otimes D_\delta,
\]
and its unit is \(k\) in weight zero.  Finite factorization makes every
displayed sum finite.  Define the coefficientwise Grothendieck group
\[
 \widehat K_{0,s}^{Q_+}(k)
 =\prod_{\beta\in Q_+}
 K_0\bigl(\operatorname{Perf}(\operatorname{sVect}_k)\bigr)
\]
and the homomorphism
\[
 \begin{aligned}
 \chi_{Q_+}^{\mathrm{tot}}:
 \widehat K_{0,s}^{Q_+}(k)&\longrightarrow
 \widehat{\mathbb Z[Q_+]},\\
 (c_\beta)_\beta&\longmapsto
 \sum_{\beta\in Q_+}\sdim(c_\beta)X^\beta.
 \end{aligned}
\]
Here \(\sdim(c_\beta)\) first takes the homological Euler class and
then subtracts intrinsic odd parity.  For a weightwise perfect
supercomplex \(C\), write
\(
 \chi_{Q_+}^{\mathrm{tot}}(C)
 :=\chi_{Q_+}^{\mathrm{tot}}([C])
\)
for the composite through its coefficientwise Grothendieck class.
Then
\[
 \begin{split}
 \chi_{Q_+}^{\mathrm{tot}}(C)
 &=
 \sum_{\beta,r}(-1)^r
 \sdim H_r(C)_\beta\,X^\beta\\
 &=
 \sum_{\beta,r,\epsilon}(-1)^{r+\epsilon}
 \dim H_r(C)_{\beta,\bar\epsilon}\,X^\beta
 \in\widehat{\mathbb Z[Q_+]}.
 \end{split}
\]

\begin{theorem}[The algebraic bar denominator]
\label{thm:algebraic-denominator}
Under the preceding hypotheses,
\(\mathcal B(\mathfrak n_+)\) is weightwise perfect.  Super
antisymmetrization gives a natural root-graded quasi-isomorphism
\[
 \operatorname{Alt}:
 C_\bullet^{\CE}(\mathfrak n_+;k)
 =
 \bigl(\Sym_{\operatorname{Ch}(\operatorname{sVect}_k)}
       (s\mathfrak n_+),d_{\CE}\bigr)
 \xrightarrow{\ \simeq\ }
 \mathcal B(\mathfrak n_+),
\]
where \(d_{\CE}\) lowers homological degree by one and preserves
intrinsic parity.  The map \(\operatorname{Alt}\) preserves homological
degree, intrinsic parity, and \(Q_+\)-weight.  Moreover,
\[
 \chi_{Q_+}^{\mathrm{tot}}\bigl(\mathcal B(\mathfrak n_+)\bigr)
 =
 \prod_{\alpha\in Q_+\setminus\{0\}}
 (1-X^\alpha)^{\sdim\mathfrak g_\alpha}.
\]
\end{theorem}

\begin{proof}
For homogeneous \(x_1,\ldots,x_r\), the comparison is
\[
 x_1\wedge\cdots\wedge x_r
 \longmapsto
 \sum_{\sigma\in S_r}
 \operatorname{sgn}(\sigma)\kappa(\sigma;x)
 [x_{\sigma(1)}|\cdots|x_{\sigma(r)}],
\]
where \(\kappa(\sigma;x)\) is the super Koszul sign.  Compatibility with
the differentials is the super antisymmetrization calculation.  Filter
both sides by PBW degree.  In each fixed \(Q_+\)-weight this filtration
is finite and exhaustive.  Super PBW gives
\[
 \operatorname{gr}U(\mathfrak n_+)
 \cong
 \Sym(\mathfrak n_{+,\bar0})
 \otimes\Lambda(\mathfrak n_{+,\bar1}).
\]
Put
\[
 S=\Sym(\mathfrak n_{+,\bar0})
   \otimes\Lambda(\mathfrak n_{+,\bar1}).
\]
The associated-graded target is \(N\Barc_\bullet(k,S,k)\), whereas the
associated-graded source has zero differential and is
\[
 \Lambda(\mathfrak n_{+,\bar0})
 \otimes\Sym(
   \mathfrak n_{+,\bar1})
 \cong
 \Sym_{\operatorname{Ch}(\operatorname{sVect}_k)}(s\mathfrak n_+).
\]
For one even generator \(x\), the Koszul resolution of \(k\) over
\(k[x]\) shows that the normalized bar has homology
\(k\oplus k[sx]\), represented by \(1\) and \([x]\).  For one odd
generator \(\theta\), the augmentation ideal of \(\Lambda(\theta)\) is
\(k\theta\) with \(\theta^2=0\), so its normalized bar differential
vanishes and
\[
 \operatorname{Alt}((s\theta)^r)
 =r![\theta|\cdots|\theta].
\]
Thus both one-generator comparisons are quasi-isomorphisms in
characteristic zero.  The shuffle Eilenberg--Zilber map and the
K\"unneth theorem identify \(\operatorname{gr}\operatorname{Alt}\), in
each fixed weight, with their finite tensor product.  Hence it is a
quasi-isomorphism.  The finite exhaustive PBW filtration in each weight
then proves that \(\operatorname{Alt}\) is a quasi-isomorphism.
The ordinary exterior resolution inside the normalized standard
resolution is
\cite{CE56}*{Chapter XIII, Section 7, Theorem 7.1}; the same filtration
argument, with the displayed Koszul signs, gives the super extension.
The PBW and chain/cochain conventions are in
\cite{CE56}*{Chapter XIII, Sections 3 and 8}.

The finite-factorization property and finite-dimensional root spaces
make every fixed weight a bounded finite-dimensional complex.
Euler--Poincaré therefore reduces the character to the underlying
superobject
\(\Sym_{\operatorname{Ch}(\operatorname{sVect}_k)}(s\mathfrak n_+)\);
the differential
\(d_{\CE}\), which is generally nonzero, does not change it.  An
intrinsically even root vector becomes total-odd after suspension and
contributes \(1-X^\alpha\).  An intrinsically odd root vector becomes
total-even and contributes \((1-X^\alpha)^{-1}\).  Multiplication over
a homogeneous basis gives the stated exponent and product.
For a one-dimensional abelian Lie superalgebra in degree \(\alpha\),
with its generator respectively even or odd, these two cases are
\[
 1-X^\alpha,\qquad
 1+X^\alpha+X^{2\alpha}+\cdots=(1-X^\alpha)^{-1},
\]
which fixes the homological and intrinsic-parity signs independently.
\end{proof}

\begin{remark}
Theorem~\ref{thm:algebraic-denominator} is an ordinary algebraic result.
It has no Ran descent datum.  Its product depends only on signed root
spaces; a generalized Cartan datum, positive system, Weyl vector, and
imaginary-simple-root correction are additional inputs to a
superdenominator identity.
\end{remark}

\begin{problem}[The carrier-changing comparison]
\label{prob:carrier-change}
For this problem set \(k=\mathbb C\); more generally, a symmetric
monoidal change-of-coefficients functor must be included among the data.
Given an ordered factorization algebra \(A\), construct a
grading-preserving comparison between its realized factorization bar and
\(\mathcal B(\mathfrak n_+)\).  One sufficient set of data is a strong
monoidal functor
\[
 \Phi:
 (\FactD(X),\otimes^!,\boldsymbol\omega_{\Ran})
 \longrightarrow
 \left(
 \widehat{\operatorname{Ch}}^{\mathrm{pf}}_{Q_+}
   (\operatorname{sVect}_{\mathbb C}),
 \widehat\otimes,\mathbb C\right),
 \qquad
 \Phi(A)\simeq U(\mathfrak n_+)
\]
as augmented \(E_1\)-algebras, together with preservation of the chosen
realization.  For a strict augmented dg representative, the comparison
must already commute with the length-two bar face induced by
multiplication.  If \(\eta:\Phi(A)\to U(\mathfrak n_+)\) denotes its
algebra map and \(\phi_2\) the tensor comparison of \(\Phi\), this equation is
\[
 \eta\,\Phi(\mu_A)\,\phi_2
 =\mu_U(\eta\otimes\eta),
 \qquad
 \varepsilon_U\eta=\Phi(\varepsilon_A),
\]
where the unit comparison identifies \(\Phi(\boldsymbol\omega_{\Ran})\)
with \(\mathbb C\).  Here \(\mu_A,\mu_U\) are the two products and
\(\varepsilon_A,\varepsilon_U\) their augmentations.
For a homotopy-coherent comparison these equalities require specified
homotopies and their higher compatibility equations.
Alternatively, a direct construction must first map
both bars to one named root-graded chain category and then give a
grading-preserving quasi-isomorphism there.
\end{problem}

\begin{conditionaltheorem}[Comparison through a monoidal functor]
\label{thm:conditional-comparison}
Assume the first set of data in
Problem~\ref{prob:carrier-change}, including the hypotheses of
Theorem~\ref{thm:realized-bar}.  Then there are natural equivalences
\[
 \Phi\!\left(
   \left|\Barc^{\mathcal F}_\bullet
    (\boldsymbol\omega_{\Ran},A,\boldsymbol\omega_{\Ran})\right|
 \right)
 \simeq
 \left|\Barc_\bullet(\mathbb C,\Phi(A),\mathbb C)\right|
 \simeq
 \mathcal B(\mathfrak n_+).
\]
Consequently,
\[
 \chi_{Q_+}^{\mathrm{tot}}\!\left(
 \Phi\!\left(
   \left|\Barc^{\mathcal F}_\bullet
    (\boldsymbol\omega_{\Ran},A,\boldsymbol\omega_{\Ran})\right|
 \right)\right)
 =
 \prod_{\alpha\in Q_+\setminus\{0\}}
 (1-X^\alpha)^{\sdim\mathfrak g_\alpha}.
\]
\end{conditionaltheorem}

\begin{proof}
Strong monoidality identifies the unit, multiplication, and augmentation
in every simplicial degree.  Preservation of realization gives the
first equivalence, and the augmented-algebra equivalence gives the
second.
\end{proof}

\section{From a bar product to a superdenominator}

Let \(H\) be a real bilinear space and let
\(G=G(H,\{h_i\}_{i\in I},\tau)\) be a Borcherds realization in the
sense of Gritsenko--Nikulin, Section~6.  Thus
\((h_i,h_j)\le0\) for \(i\ne j\),
\(2(h_i,h_j)/(h_i,h_i)\) is integral whenever \((h_i,h_i)>0\), and
\((h_i,h_i)\le0\) for every odd simple index \(i\in\tau\).  Assume the
hypotheses of their Statement~6.8\(^{\prime}\), so in particular there
are no odd real simple roots.  Put
\[
 \widetilde Q=\bigoplus_{i\in I}\mathbb Zr_i,\qquad
 \widetilde Q_+=\sum_{i\in I}\mathbb N r_i,\qquad
 \varpi(r_i)=h_i,
\]
and require the conditions (6.29)--(6.31): every fibre of
\(\varpi|_{\widetilde Q_+}\) is finite and
\[
 \varpi^{-1}(0)\cap\widetilde Q_+=\{0\}.
\]
Let \(Q=\sum_i\mathbb Zh_i\subset H\),
\(Q_+=\varpi(\widetilde Q_+)\), and
\(\Delta_+=\Delta\cap(Q_+\setminus\{0\})\).  These conditions make
\(Q_+\) pointed with the finite-factorization property: every ordered
decomposition lifts to an ordered decomposition of one of the finitely
many elements in its \(\varpi\)-fibre, and an element of the free monoid
\(\widetilde Q_+\) has only finitely many such decompositions.  Roots
\(\alpha\in\Delta_+\) below are the \(Q\)-valued roots, and
\(G_\alpha\) denotes the total finite-dimensional root space in that
degree.  Set
\[
 G_+=\bigoplus_{\alpha\in\Delta_+}G_\alpha,\qquad
 \operatorname{mult}\alpha
 =\dim G_{\alpha,\bar0}-\dim G_{\alpha,\bar1}.
\]

Let \(P\supset Q\) be a \(W\)-stable weight lattice containing a Weyl
vector \(\rho_{\mathrm{alg}}\) with
\[
 (\rho_{\mathrm{alg}},h_i)=\frac12(h_i,h_i).
\]
Here \(W\) is the Weyl group of the real simple roots and
\(\varepsilon_W(w)=(-1)^{\ell(w)}\).  Write
\(e(\lambda)e(\mu)=e(\lambda+\mu)\) and
\(X^\alpha=e(-\alpha)\), and define the shifted completion
\[
 e(\rho_{\mathrm{alg}})\widehat{\mathbb Z[Q_+]}
 :=
 \prod_{\beta\in Q_+}
 \mathbb Z\,e(\rho_{\mathrm{alg}}-\beta).
\]

\begin{corollary}[Formal superdenominator comparison]
\label{cor:formal-superdenominator}
For
\[
 \mathcal B(G_+)
 =N\Barc_\bullet\bigl(\mathbb C,U(G_+),\mathbb C\bigr),
\]
the unshifted bar character is
\[
 \chi_{Q_+}^{\mathrm{tot}}\bigl(\mathcal B(G_+)\bigr)
 =
 \prod_{\alpha\in\Delta_+}
 (1-X^\alpha)^{\operatorname{mult}\alpha}.
\]
Its Weyl shift satisfies, coefficientwise in the shifted completion,
\[
 e(\rho_{\mathrm{alg}})
 \chi_{Q_+}^{\mathrm{tot}}\bigl(\mathcal B(G_+)\bigr)
 =
 e(\rho_{\mathrm{alg}})
 \prod_{\alpha\in\Delta_+}
 (1-e(-\alpha))^{\operatorname{mult}\alpha}.
\]
It is the product side of
\[
 e(\rho_{\mathrm{alg}})
 \prod_{\alpha\in\Delta_+}
 (1-e(-\alpha))^{\operatorname{mult}\alpha}
 =
 \sum_{w\in W}\varepsilon_W(w)\,w(S),
\]
where
\[
 S=e(\rho_{\mathrm{alg}})
   \sum_s\varepsilon_I(s)e(-\varpi(s)).
\]
Here \(s=0\) is included, with \(\varepsilon_I(0)=1\).  The remaining
\(s\) range over sums of pairwise perpendicular imaginary simple
indices.  If \(s\) contains \(n_0\) even and \(n_1\) odd imaginary
simple indices, then \(\varepsilon_I(s)=(-1)^{n_0}\); indices outside
the set of odd isotropic simple roots are required to be distinct, and
\(\varepsilon_I(s)=0\) for an inadmissible sum.
\end{corollary}

\begin{proof}
Theorem~\ref{thm:algebraic-denominator} gives the product on the left.
The remaining identity is the superdenominator formula
\cite{GN97}*{the paragraph before Statement 6.5 and Statement
6.8\(^{\prime}\), with (6.29)--(6.31)}.
\end{proof}

\begin{remark}
For ordinary generalized Kac--Moody algebras the corresponding formula
is \cite{Bor88}*{formula (3), pp.~506--507}.  That formula does not
contain the super sign or the odd-imaginary-root correction above.
\end{remark}

\section{The Igusa specialization}

\subsection{The active root cone}

Let
\[
 M_0=\mathbb Zf_2\oplus\mathbb Zf_3\oplus\mathbb Zf_{-2},
 \qquad
 (f_2,f_{-2})=-1,\qquad (f_3,f_3)=2,
\]
with all other basis pairings zero, and put
\[
 \delta_1=2f_2-f_3,\qquad
 \delta_2=2f_{-2}-f_3,\qquad
 \delta_3=f_3.
\]
Then
\[
 Q=M_{II}
 =\mathbb Z\delta_1\oplus\mathbb Z\delta_2
  \oplus\mathbb Z\delta_3,\qquad
 Q_+=\mathbb N\delta_1+\mathbb N\delta_2+\mathbb N\delta_3.
\]
Thus \(Q_+\cong\mathbb N^3\) has the finite-factorization property.
These lattice and chamber data are those of
\cite{GN97}*{(2.1)--(2.4)}.

For \(n,\ell,m\in\mathbb Z\), define
\[
 \alpha(n,\ell,m)=2nf_2-\ell f_3+2mf_{-2}.
\]
Its coordinates in \(Q\) are
\begin{equation}
 \alpha(n,\ell,m)
 =n\delta_1+m\delta_2+(n+m-\ell)\delta_3.
 \label{eq:igusa-root-coordinate}
\end{equation}
The positive triples satisfy \(n,m\ge0\), with arbitrary \(\ell\) if
\(n>0\) or \(m>0\), and with \(\ell<0\) if \(n=m=0\)
\cite{GN97}*{Theorem 4.1 and Remark 4.2}.

Let
\[
 \phi_{0,1}(\tau,z)
 =\sum_{N\ge0,\ \ell\in\mathbb Z}
 f(N,\ell)e^{2\pi i(N\tau+\ell z)}
\]
be the weak Jacobi form normalized by
\(f(0,0)=10\) and \(f(0,\pm1)=1\).  It has weight zero and index one,
and its support satisfies
\[
 f(N,\ell)\ne0\quad\Longrightarrow\quad
 4N-\ell^2\ge-1.
\]
These assertions follow from
\cite{GN97}*{(4.11) and the paragraph following it}.

\begin{lemma}[Active positive cone]
\label{lem:igusa-active-cone}
If \((n,\ell,m)>0\) and \(f(nm,\ell)\ne0\), then
\(\alpha(n,\ell,m)\in Q_+\).
\end{lemma}

\begin{proof}
If \(n+m>0\) and \(\ell\ge n+m+1\), then
\[
 \ell^2\ge(n+m+1)^2>4nm+1,
\]
contrary to the support inequality.  Hence \(n+m-\ell\ge0\).
If \(n=m=0\), positivity and the support inequality force
\(\ell=-1\).  Equation~\eqref{eq:igusa-root-coordinate} gives the
result.
\end{proof}

The restriction to active roots is essential.  The raw positive triples
have infinitely many decompositions
\[
 (1,0,0)=(1,j,0)+(0,-j,0),\qquad j\ge1.
\]
Both summands are simultaneously active only for \(j=1\).  At the
boundary \(n=m=0\), the cone selects
\(\delta_3=\alpha(0,-1,0)\).  Under the analytic substitution below its
factor is \(1-y^{-2}\); reversing the root changes it by the exact
monomial factor
\[
 1-y^2=-y^2(1-y^{-2}).
\]

Gritsenko and Nikulin construct a generalized Kac--Moody
superalgebra \(\mathfrak g_{\Delta_5}\) with finite-dimensional root
spaces and
\[
 \sdim(\mathfrak g_{\Delta_5})_{\alpha(n,\ell,m)}
 =f(nm,\ell)
\]
\cite{GN97}*{Section 3 and (3.4)}.  The finite-fibre conditions ensuring
finite-dimensional projected root spaces are
\cite{GN97}*{(6.29)--(6.31)}.  Define
\[
 \mathfrak n_{\Delta_5,+}
 =\bigoplus_{\alpha\in\Delta_+}
   (\mathfrak g_{\Delta_5})_\alpha,\qquad
 \mathcal B_{\Delta_5}
 =N\Barc_\bullet\bigl(
   \mathbb C,U(\mathfrak n_{\Delta_5,+}),\mathbb C\bigr).
\]
By Lemma~\ref{lem:igusa-active-cone}, this is weightwise perfect in
\(\operatorname{Ch}_{\ge0}
(\operatorname{sVect}_{\mathbb C}^{Q_+})\).

\subsection{The Weyl shift and the formal Fourier map}

The geometric Weyl vector is
\[
 \rho_{\mathrm{geo}}
 =\frac12(\delta_1+\delta_2+\delta_3),\qquad
 (\rho_{\mathrm{geo}},\delta_i)=-1.
\]
The algebraic convention in the superdenominator formula is
\[
 \rho_{\mathrm{alg}}=-\rho_{\mathrm{geo}},\qquad
 (\rho_{\mathrm{alg}},\delta_i)=1.
\]
These conventions are respectively
\cite{GN97}*{(2.3)--(2.4), (3.2)--(3.3)} and the paragraph preceding
\cite{GN97}*{Statement 6.5, together with Statement
6.8\(^{\prime}\)}.  In particular,
\[
 e(\rho_{\mathrm{alg}})=e(-\rho_{\mathrm{geo}})
 =X^{\rho_{\mathrm{geo}}},\qquad X^\alpha:=e(-\alpha).
\]
Since \(\rho_{\mathrm{geo}}\) lies in the weight lattice but not in
\(Q\), this expression belongs to the shifted module
\[
 X^{\rho_{\mathrm{geo}}}\widehat{\mathbb Z[Q_+]}.
\]

Let
\[
 \widehat{\mathscr F}_{Q_+}=\mathbb Z[[u,v,w]]
\]
with its \((u,v,w)\)-adic topology.  Let
\(\omega\widehat{\mathscr F}_{Q_+}\) be the free rank-one shifted
module obtained in the enlarged lattice algebra by imposing
\(\omega^2=uvw\).  There is a continuous formal isomorphism
\[
 \mathcal F_2:\widehat{\mathbb Z[Q_+]}
 \longrightarrow\widehat{\mathscr F}_{Q_+},
 \qquad
 X^{\delta_1}\mapsto u,\quad
 X^{\delta_2}\mapsto v,\quad
 X^{\delta_3}\mapsto w,
\]
and a shifted extension
\[
 \mathcal F_2^\rho:
 X^{\rho_{\mathrm{geo}}}\widehat{\mathbb Z[Q_+]}
 \longrightarrow\omega\widehat{\mathscr F}_{Q_+},
 \qquad
 \mathcal F_2^\rho(X^{\rho_{\mathrm{geo}}}F)
 =\omega\,\mathcal F_2(F).
\]

\begin{proposition}[Formal Igusa character]
\label{prop:formal-igusa-character}
One has
\[
 \mathcal F_2^\rho\!\left(
 X^{\rho_{\mathrm{geo}}}
 \chi_{Q_+}^{\mathrm{tot}}(\mathcal B_{\Delta_5})\right)
 =
 \omega\prod_{\substack{(n,\ell,m)>0\\ f(nm,\ell)\ne0}}
 \left(1-u^nv^mw^{\,n+m-\ell}\right)^{f(nm,\ell)}
\]
in \(\omega\widehat{\mathscr F}_{Q_+}\).
\end{proposition}

\begin{proof}
Apply Theorem~\ref{thm:algebraic-denominator} and the signed
multiplicity formula (3.4) of \cite{GN97}.  Equation
\eqref{eq:igusa-root-coordinate} gives the exponent of
\(u^nv^mw^{n+m-\ell}\), while the Weyl-vector convention gives the
factor \(\omega\).
\end{proof}

\begin{remark}
Proposition~\ref{prop:formal-igusa-character} is a formal equality in a
shifted cone module.  Neither the scalar \(64\) nor analytic convergence
can be recovered from the exponent table.
\end{remark}

\subsection{Analytic realization at the cusp}

For
\[
 Z=\begin{pmatrix}z_1&z_2\\z_2&z_3\end{pmatrix}\in\mathbb H_2,
 \qquad
 q=e^{2\pi iz_1},\quad y=e^{2\pi iz_2},\quad
 p=e^{2\pi iz_3},
\]
set
\[
 u_Z=q^2y^2,\qquad v_Z=p^2y^2,\qquad
 w_Z=y^{-2},\qquad \omega_Z=qyp.
\]
Then
\[
 \omega_Z^2=u_Zv_Zw_Z,\qquad
 u_Z^nv_Z^mw_Z^{\,n+m-\ell}
 =q^{2n}y^{2\ell}p^{2m}.
\]

Writing \(y_j=\operatorname{Im}z_j\), define
\[
 U_T=\left\{Z\in\mathbb H_2:
 -y_2>T,\quad y_1+y_2>T,\quad y_3+y_2>T\right\}.
\]
For
\(\zeta=z_3f_2+z_2f_3+z_1f_{-2}\), these inequalities are equivalent to
\((\operatorname{Im}\zeta,\delta_i)<-2T\).  In particular,
\[
 |u_Z|,|v_Z|,|w_Z|<e^{-4\pi T}.
\]

Let
\(\widehat{\mathscr F}_{Q_+}^{\mathrm{abs}}(U_T)\) consist of the
series \(F=\sum A_{a,b,c}u^av^bw^c\) for which
\[
 \|F\|_K
 =\sum_{a,b,c\ge0}|A_{a,b,c}|
 \sup_{Z\in K}|u_Z^av_Z^bw_Z^c|<\infty
\]
for every compact \(K\Subset U_T\).  Endow this algebra with the
seminorms \(\|\cdot\|_K\), and \(\mathcal O(U_T)\) with the
compact-open topology.  Termwise summation defines continuous maps
\[
 \operatorname{an}_{U_T}:
 \widehat{\mathscr F}_{Q_+}^{\mathrm{abs}}(U_T)
 \longrightarrow\mathcal O(U_T),
\]
\[
 \operatorname{an}_{U_T}^{\rho}:
 \omega\widehat{\mathscr F}_{Q_+}^{\mathrm{abs}}(U_T)
 \longrightarrow\mathcal O(U_T),\qquad
 \omega F\longmapsto(qyp)\operatorname{an}_{U_T}(F).
\]
There is no such termwise realization on the whole adic completion:
\(\sum_{r\ge0}r!u^r\) has radius of convergence zero.

For an active triple put
\[
 d=n+m+(n+m-\ell).
\]
The support inequality gives
\[
 \sqrt{\max(0,4nm-\ell^2)}\le n+m\le d,
\]
and only \(O(d^2)\) triples have a fixed value of \(d\).  The coefficient
estimate in the proof of \cite{GN97}*{Theorem 4.1}, together with
\[
 \left|u_Z^nv_Z^mw_Z^{n+m-\ell}\right|
 \le e^{-4\pi Td},
\]
therefore implies that, for sufficiently large \(T\),
\[
 \sum_{\substack{(n,\ell,m)>0\\ f(nm,\ell)\ne0}}
 |f(nm,\ell)|\sup_{Z\in K}
 \frac{|q^{2n}y^{2\ell}p^{2m}|}
      {1-|q^{2n}y^{2\ell}p^{2m}|}<\infty
\]
for every compact \(K\Subset U_T\).  Consequently the formal product
in Proposition~\ref{prop:formal-igusa-character} belongs to
\(\omega\widehat{\mathscr F}_{Q_+}^{\mathrm{abs}}(U_T)\): for a factor
with exponent \(a\), its absolute binomial norm is bounded by
\(\exp(|a|r/(1-r))\) when the monomial has modulus at most \(r<1\).

Let \(\Delta_5\) be the product of the ten even genus-two theta
constants, and put \(D_5=64^{-1}\Delta_5\).

\begin{theorem}[Igusa realization]
\label{thm:igusa-realization}
For sufficiently large \(T\),
\[
\begin{aligned}
 &\operatorname{an}_{U_T}^{\rho}\!\left(
 \mathcal F_2^\rho\!\left(
 X^{\rho_{\mathrm{geo}}}
 \chi_{Q_+}^{\mathrm{tot}}(\mathcal B_{\Delta_5})\right)\right)\\
 &\qquad =
 qyp\prod_{\substack{(n,\ell,m)>0\\ f(nm,\ell)\ne0}}
 \left(1-q^{2n}y^{2\ell}p^{2m}\right)^{f(nm,\ell)}\\
 &\qquad =D_5(2Z)=64^{-1}\Delta_5(2Z)
\end{aligned}
\]
as holomorphic functions on \(U_T\).
\end{theorem}

\begin{proof}
The product formula and its convergence are
\cite{GN97}*{Theorem 4.1 and its proof, Remark 4.2}, evaluated at
\(2Z\).  The normalization \(D_5=64^{-1}\Delta_5\) follows from the
ten-theta product and its leading Fourier coefficient
\cite{GN97}*{(1.2), (1.6)--(1.7)}.  The identification with the
superdenominator is \cite{GN97}*{Proposition 3.1}; the scale
\(Z\mapsto2Z\) is explained in \cite{GN97}*{Remark 3.2}.
\end{proof}

\begin{remark}[Automorphic descent]
\label{rem:automorphic-descent}
The form \(\Delta_5\) satisfies
\[
 \Delta_5(M\langle Z\rangle)
 =v(M)\det(CZ+D)^5\Delta_5(Z),
 \qquad
 M=\begin{pmatrix}A&B\\C&D\end{pmatrix}
 \in\operatorname{Sp}_4(\mathbb Z),
\]
where \(v\) is the nontrivial multiplier of
\cite{GN97}*{(1.3)--(1.5)}.  Thus it is a section of the weight-five
automorphic line with multiplier \(v\).

Let
\[
 g_2=\begin{pmatrix}2I_2&0\\0&I_2\end{pmatrix}
 \in\operatorname{GSp}_4(\mathbb Q),\qquad
 \Gamma^{(2)}=g_2^{-1}\operatorname{Sp}_4(\mathbb Z)g_2.
\]
For \(\gamma\in\Gamma^{(2)}\), put
\(M=g_2\gamma g_2^{-1}\).  Then
\[
 D_5\!\left(2\,\gamma\langle Z\rangle\right)
 =
 v(M)\det\!\left(C_M(2Z)+D_M\right)^5D_5(2Z).
\]
Accordingly, \(D_5(2Z)\) is the pullback of the weight-five automorphic
section along the scaled cusp embedding.  The formal bar character has
no equivariance datum of this kind; automorphic descent is supplied by
the separate Gritsenko--Nikulin theorem.
\end{remark}

\section{The positive Witt algebra}

Let
\[
 L_1^{\mathrm{poly}}=\bigoplus_{j\ge1}\mathbb Ce_j,
 \qquad
 \widehat L_1=\prod_{j\ge1}\mathbb Ce_j,
 \qquad [e_i,e_j]=(j-i)e_{i+j}.
\]
Equip \(\widehat L_1\) with the filtration topology
\[
 F^N\widehat L_1=\prod_{j\ge N}\mathbb Ce_j.
\]
For every \(n,w\), evaluation gives a perfect pairing
\[
 C_n^{\CE}(L_1^{\mathrm{poly}};\mathbb C)_w
 \otimes
 C^n_{\mathrm{cont}}(\widehat L_1;\mathbb C)_w
 \longrightarrow\mathbb C.
\]
Indeed, only finitely many strictly increasing positive indices sum to
\(w\).  The chain and cochain differentials are adjoint, hence
\begin{equation}
 H_n^{\CE}(L_1^{\mathrm{poly}};\mathbb C)_w
 \cong
 H^n_{\mathrm{cont}}(\widehat L_1;\mathbb C)_w^\vee.
 \label{eq:witt-duality}
\end{equation}
The chain--cochain pairing is the one in
\cite{Gon73}*{equations (2)--(3)}.
The bar and CE chain constructions are covariant; continuous cochains
are contravariant.  The duality in \eqref{eq:witt-duality} is the
variance-reversing bridge.

\begin{theorem}[Positive Witt bar homology]
\label{thm:witt-homology}
For \(n\ge1\),
\[
 \dim H_n\mathcal B(L_1^{\mathrm{poly}})_w
 =
 \begin{cases}
 1,&w=(3n^2-n)/2\ \text{or}\ w=(3n^2+n)/2,\\
 0,&\text{otherwise}.
 \end{cases}
\]
Thus
\(\dim H_n\mathcal B(L_1^{\mathrm{poly}})=2\) for every \(n\ge1\).
\end{theorem}

\begin{proof}
Theorem~\ref{thm:algebraic-denominator} identifies the normalized
enveloping-algebra bar with the CE chain complex, weight by weight.
Equation~\eqref{eq:witt-duality} identifies its homology with the dual
of continuous cohomology.  Goncharova's Main Theorem and Example 1.1
give precisely the two displayed weights and the vanishing elsewhere
\cite{Gon73}*{Sections 1.1--1.3, pp.~8--10}.
\end{proof}

The smallest weight calculation is
\[
 [L_1^{\mathrm{poly}},L_1^{\mathrm{poly}}]
 =\operatorname{span}\{e_j:j\ge3\},
\]
so \(H_1\) has weights \(1\) and \(2\).

\begin{corollary}[Pentagonal character]
\label{cor:pentagonal-character}
In \(\mathbb Z[[q]]\),
\[
\begin{aligned}
 \chi_q^{\mathrm{tot}}\mathcal B(L_1^{\mathrm{poly}})
 &=\prod_{m\ge1}(1-q^m)\\
 &=1+\sum_{n\ge1}(-1)^n
 \left(q^{(3n^2-n)/2}+q^{(3n^2+n)/2}\right).
\end{aligned}
\]
\end{corollary}

\begin{proof}
The first equality is
Theorem~\ref{thm:algebraic-denominator}.  The second is Euler's
pentagonal number theorem
\cite{DLMF}*{Section 27.14(ii), equations (27.14.4)--(27.14.5)}.
\end{proof}

\begin{remark}
The product corroborates the weights in
Theorem~\ref{thm:witt-homology}; it does not prove the vanishing
statement.  Euler characteristic cannot see a direct sum of
one-dimensional zero-differential chain complexes in two adjacent
degrees at the same weight.  The absence of additional classes comes
from Goncharova's theorem.
\end{remark}

\section{Motzkin--Riordan identities and ordinary carriers}

Let \(M_n\) count Motzkin paths of length \(n\), and let \(R_n\) count
those with no horizontal step at height zero.  Put
\[
 M_0=R_0=1,\qquad R_1=0,
\]
and define
\[
 M(z)=\sum_{n\ge0}M_nz^n,\qquad
 R(z)=\sum_{n\ge0}R_nz^n.
\]

\begin{proposition}[The Motzkin--Riordan relation]
\label{prop:motzkin-riordan}
The generating functions satisfy
\[
 M=1+zM+z^2M^2,\qquad R=1+z^2MR.
\]
Consequently,
\[
 R=\frac{M}{1+zM},\qquad
 M=R+\frac{R-1}{z},
\]
and, for every \(n\ge0\),
\[
 M_n=R_n+R_{n+1},\qquad
 M_{n+1}-M_n=R_{n+2}-R_n.
\]
Moreover,
\[
 M(z)=\frac{1-z-\sqrt{1-2z-3z^2}}{2z^2},
\qquad
 R(z)=\frac{1+z-\sqrt{1-2z-3z^2}}{2z(1+z)}.
\]
\end{proposition}

\begin{proof}
A nonempty Motzkin path is either a horizontal step followed by a
Motzkin path, or has a unique first-return decomposition \(UPDQ\), with
\(P,Q\) Motzkin paths.  This gives the first equation for \(M\).  A
nonempty Riordan path has the same form \(UPDQ\), now with \(P\)
Motzkin and \(Q\) Riordan, which gives the equation for \(R\).
Algebraic elimination yields the two rational relations and coefficient
comparison gives the identities.  Solving the quadratic equation for
the branch with constant term \(1\) gives the displayed closed forms.
These identities and generating equations also appear in
\cite{Ber99}*{pp.~87, 89, 100--101}; the path convention is
\cite{CDY08}*{Section 2}.
\end{proof}

\begin{remark}
Splitting a Motzkin path according to whether its first step is
horizontal does not prove \(M_n=R_n+R_{n+1}\).  At length three the
paths
\[
 HHH,\qquad HUD,\qquad UDH,\qquad UHD
\]
split as \(2+2\), whereas \(R_3+R_4=1+3\).  The common square-root
discriminant follows from Proposition~\ref{prop:motzkin-riordan}.  This
scalar identity neither constructs nor rules out a
representation-theoretic or Drinfeld--Sokolov comparison.
\end{remark}

\begin{proposition}[The ordinary \(\mathfrak{sl}_2\) calculation]
\label{prop:sl2}
Let \(k\) have characteristic zero and give
\(\mathfrak{sl}_2(k)\) the trivial module \(k\).  Then
\[
 H^\bullet(\mathfrak{sl}_2;k)\cong\Lambda(c_3),
 \qquad |c_3|=3,
\]
and the CE cochain dimensions are \((1,3,3,1)\).  In contrast, every
positive simplicial degree of
\[
 N\Barc_\bullet(k,U(\mathfrak{sl}_2),k)
\]
is infinite-dimensional.
\end{proposition}

\begin{proof}
Use
\([h,e]=2e,\ [h,f]=-2f,\ [e,f]=h\).  On the dual basis,
\[
 de^\vee=2e^\vee\wedge h^\vee,\qquad
 df^\vee=-2f^\vee\wedge h^\vee,\qquad
 dh^\vee=-e^\vee\wedge f^\vee.
\]
Thus \(d:C^1\to C^2\) is an isomorphism and
\(d:C^2\to C^3\) is zero.  The class
\[
 c_3=[e^\vee\wedge f^\vee\wedge h^\vee]
\]
generates \(H^3\), which proves the cohomology statement directly.

For the normalized associative bar,
\[
 \bigl(N\Barc_\bullet(k,U(\mathfrak{sl}_2),k)\bigr)_r
 \cong\overline{U(\mathfrak{sl}_2)}^{\otimes r}.
\]
PBW gives linearly independent elements
\(e,e^2,e^3,\ldots\) in the augmentation ideal.  Length one is already
infinite-dimensional, and hence so is every positive length.
\end{proof}

\begin{remark}[Ordinary and chiral carriers]
\label{rem:family-carriers}
For the standard Virasoro bracket,
\(e_n\mapsto-L_n\) identifies \(L_1^{\mathrm{poly}}\) with the
positive-mode Lie subalgebra, because the central cocycle vanishes on
positive indices.  This inclusion does not identify the ordinary
enveloping-algebra bar with a chiral bar of a Virasoro vacuum object.
Such a comparison must specify the central charge, vacuum object,
completion, homological and conformal gradings, and a map relating the
two complexes.

Likewise, Proposition~\ref{prop:sl2} has no implication for an affine
or chiral \(\widehat{\mathfrak{sl}}_2\) object without a specified
level, vacuum module, completion, grading, and comparison.  A
grading-preserving quasi-isomorphism would suffice.  A spectral-sequence
argument additionally requires an identified abutment, strong
convergence, and control of all relevant differentials.
\end{remark}

For the ordinary carriers, consider the two proposed scalar assignments
\[
 \dim H_n\mathcal B(L_1^{\mathrm{poly}})=M_{n+1}-M_n,
 \qquad
 \dim H^n(\mathfrak{sl}_2;k)=R_{n+3}.
\]
Their first discrepancies are decisive:
\[
 M_2-M_1=1\ne2=\dim H_1\mathcal B(L_1^{\mathrm{poly}}),
\qquad
 R_4=3\ne0=\dim H^1(\mathfrak{sl}_2;k).
\]
A finite total Betti number in each degree is a meaningful invariant;
the issue is not its scalar type but the absence of a comparison with
the proposed chiral or affine carrier.

\section{From operations to the cusp product}

The algebraic product follows from a chain comparison and finite
weightwise Euler characteristics.  Its analytic realization then uses
the specified Fourier map and cusp domain.  The following functors make
these successive operations explicit and identify where a geometric
factorization comparison would enter.

Fix a finite-factorization monoid \(Q_+\).  Let
\(\operatorname{LieSup}^{Q_+,\mathrm{fd},+}_k\) denote the category of
positively \(Q_+\)-graded Lie superalgebras with finite-dimensional
pieces.  Let
\(\operatorname{Alg}^{\mathrm{aug},Q_+,\mathrm{fd},+}_{E_1}\)
denote augmented \(Q_+\)-graded superalgebras \(B\) with \(B_0=k\),
augmentation ideal supported away from zero, and finite-dimensional
homogeneous pieces.  Enveloping algebra and normalized bar are functors
\[
\begin{tikzcd}[column sep=3.2em]
 \operatorname{LieSup}^{Q_+,\mathrm{fd},+}_k
   \arrow[r,"U"]
 & \operatorname{Alg}^{\mathrm{aug},Q_+,\mathrm{fd},+}_{E_1}
   (\operatorname{sVect}_k)
   \arrow[r,"{N\Barc_\bullet(k,-,k)}"]
 & \widehat{\operatorname{Ch}}^{\mathrm{pf}}_{Q_+}
   (\operatorname{sVect}_k).
\end{tikzcd}
\]
Finite factorization makes the last functor land in its stated target.
Theorem~\ref{thm:algebraic-denominator} supplies the natural
bar--CE quasi-isomorphism on the first category.

For the geometric branch, the typed arrows are
\[
\begin{tikzcd}[column sep=2.5em]
 \Alg^{\mathrm{aug}}_{E_1}(\FactD(X),\otimes^!)
   \arrow[r,"\Barc^{\mathcal F}_\bullet"]
 & \operatorname{Fun}(\Delta^{op},\FactD(X))
   \arrow[r,dashed,"\lvert-\rvert"]
 & \FactD(X)
   \arrow[r,dashed,"\Phi"]
 & \widehat{\operatorname{Ch}}^{\mathrm{pf}}_{Q_+}
   (\operatorname{sVect}_{\mathbb C}).
\end{tikzcd}
\]
The upper realization arrow is conditional on
Theorem~\ref{thm:realized-bar}; \(\Phi\) is the open construction in
Problem~\ref{prob:carrier-change}.  Neither dashed arrow is used in the
algebraic theorem.

For an object \(C\) of the completed chain category, passage to its
coefficientwise class and then to its total character is the typed
sequence
\[
 \operatorname{Obj}\!\left(
   \widehat{\operatorname{Ch}}^{\mathrm{pf}}_{Q_+}
   (\operatorname{sVect}_k)\right)\big/\simeq
 \xrightarrow{[-]}
 \widehat K_{0,s}^{Q_+}(k)
 \xrightarrow{\chi_{Q_+}^{\mathrm{tot}}}
 \widehat{\mathbb Z[Q_+]}.
\]
The second arrow is a group homomorphism, not a chain functor.  The
Igusa specialization continues from the final ring by the shifted
formal map \(\mathcal F_2^\rho\), then by the restricted analytic
realization on \(U_T\).  Neither map constructs \(\Phi\).

\begin{thebibliography}{99}

\bibitem{AM10}
M.~Aguiar and S.~Mahajan,
\emph{Monoidal Functors, Species and Hopf Algebras},
CRM Monograph Series 29, American Mathematical Society, Providence,
RI, 2010.

\bibitem{Ber99}
F.\,R.~Bernhart,
\emph{Catalan, Motzkin, and Riordan numbers},
Discrete Math. \textbf{204} (1999), 73--112.

\bibitem{Bor88}
R.\,E.~Borcherds,
\emph{Generalized Kac--Moody algebras},
J. Algebra \textbf{115} (1988), 501--512.

\bibitem{CE56}
H.~Cartan and S.~Eilenberg,
\emph{Homological Algebra},
Princeton University Press, Princeton, NJ, 1956.

\bibitem{CDY08}
W.\,Y.\,C.~Chen, E.\,Y.\,P.~Deng, and L.\,L.\,M.~Yang,
\emph{Riordan paths and derangements},
Discrete Math. \textbf{308} (2008), 2222--2227.

\bibitem{DLMF}
NIST Digital Library of Mathematical Functions,
\emph{Chapter 27: Functions of number theory},
Section 27.14, \url{https://dlmf.nist.gov/27.14}.

\bibitem{FG12}
J.~Francis and D.~Gaitsgory,
\emph{Chiral Koszul duality},
Selecta Math. (N.S.) \textbf{18} (2012), 27--87.

\bibitem{Gon73}
L.\,V.~Goncharova,
\emph{The cohomologies of Lie algebras of formal vector fields on the
line},
Funktsional. Anal. i Prilozhen. \textbf{7} (1973), no.~2, 6--14;
English transl., Funct. Anal. Appl. \textbf{7} (1973), 91--97.

\bibitem{GN97}
V.\,A.~Gritsenko and V.\,V.~Nikulin,
\emph{Siegel automorphic form corrections of some Lorentzian
Kac--Moody Lie algebras},
Amer. J. Math. \textbf{119} (1997), 181--224.

\bibitem{Tor25}
T.~Torii,
\emph{On duoidal \(\infty\)-categories},
J. Homotopy Relat. Struct. \textbf{20} (2025), no.~1, 125--162.

\end{thebibliography}

\end{document}
